\section{Graph RAG versus Vector RAG}
\label{sec:gr-graphrag-vs-vectorrag}

Graph RAG is the family of RAG methods in which the external knowledge is
organized as a graph and retrieval exploits that structure, rather than treating
the corpus as a flat set of independent chunks~\cite{peng2024graphrag}. By
retrieving connected entities and the relations between them, Graph RAG captures
relational knowledge and supports more precise, context-aware
responses~\cite{peng2024graphrag}. The two paradigms are therefore complementary:
vector RAG excels at locating locally relevant passages, whereas Graph RAG adds
the relational and global structure that vector similarity
discards~\cite{peng2024graphrag, edge2024graphrag}, as contrasted in
\Cref{fig:gr-vector-vs-graph}.

\begin{figure}[htbp]
  \centering
  \begin{tikzpicture}[
    chunk/.style={draw, rounded corners, minimum width=1.5cm, minimum height=0.5cm,
                  fill=black!5, font=\scriptsize},
    ent/.style={draw, circle, minimum size=0.7cm, font=\scriptsize, inner sep=0pt},
    rel/.style={-{Latex[length=1.5mm]}, semithick},
    title/.style={font=\small\bfseries}]
    % --- left: vector RAG (flat chunks) ---
    \begin{scope}
      \node[title] (t1) at (0,2.1) {Vector RAG};
      \foreach \i in {0,1,2} \foreach \j in {0,1}
        \node[chunk] at (\j*1.7, -\i*0.75) {};
      \node[font=\scriptsize, text width=3.2cm, align=center] at (0.85,-2.6)
        {independent chunks ranked by similarity};
    \end{scope}
    % --- right: graph RAG (connected entities) ---
    \begin{scope}[xshift=6.2cm]
      \node[title] (t2) at (0.6,2.1) {Graph RAG};
      \node[ent] (a) at (0,1)     {A};
      \node[ent] (b) at (1.4,1.4) {B};
      \node[ent] (c) at (1.8,0)   {C};
      \node[ent] (d) at (0.4,-0.6){D};
      \node[ent] (e) at (2.6,0.9) {E};
      \draw[rel] (a)--(b); \draw[rel] (b)--(c); \draw[rel] (a)--(d);
      \draw[rel] (c)--(d); \draw[rel] (b)--(e); \draw[rel] (c)--(e);
      \node[font=\scriptsize, text width=3.2cm, align=center] at (1.1,-2.6)
        {connected entities and relations};
    \end{scope}
  \end{tikzpicture}
  \caption[Vector RAG versus Graph RAG]{Vector RAG treats the corpus as independent chunks ranked by
    similarity, whereas Graph RAG retrieves connected entities and the relations
    between them, preserving relational structure~\cite{peng2024graphrag}.}
  \label{fig:gr-vector-vs-graph}
\end{figure}

The practical benefits have been demonstrated from both ends of the design space.
At the ``global'' end, building an entity knowledge graph and pre-computing
community summaries lets the system answer corpus-wide sensemaking questions that
naive RAG cannot, improving the comprehensiveness and diversity of
answers~\cite{edge2024graphrag}. At the efficiency end, incorporating graph
structure into indexing and using a dual-level retrieval scheme has been shown to
improve both accuracy and response time relative to conventional
RAG~\cite{guo2024lightrag}. More broadly, unifying language models with knowledge
graphs is argued to combine the fluent generation of the former with the
explicit, verifiable factual knowledge of the latter~\cite{pan2024unifying}.
